\section*{Chapter \ref{ch:b3:point-groups} --- Symmetry and Point Groups}

\begin{solution}{exo:b3:point-groups:1}
Ethene: three perpendicular $C_2$, three planes, $i$: $D_{2h}$. \ce{XeF4} (square
planar): $C_4$, four $C_2 \perp C_4$, $\sigma_h$, $2\sigma_v$, $2\sigma_d$, $i$, $S_4$:
$D_{4h}$. \ce{PCl5} (trigonal bipyramid): $C_3$, three $C_2$, $\sigma_h$, $3\sigma_v$,
$S_3$: $D_{3h}$. 1,3,5-Trichlorobenzene: the same elements as \ce{BF3}: $D_{3h}$.
\end{solution}

\begin{solution}{exo:b3:point-groups:2}
Staggered ethane: $C_3$, three $C_2$ perpendicular, three $\sigma_d$, $i$, $S_6$:
$D_{3d}$. Eclipsed: $C_3$, three $C_2$, $\sigma_h$, three $\sigma_v$: $D_{3h}$.
Chair cyclohexane: $D_{3d}$.
\end{solution}

\begin{solution}{exo:b3:point-groups:3}
\ce{SO3} ($D_{3h}$): no. \ce{SO2} ($C_{2v}$): yes. \ce{PCl3} ($C_{3v}$): yes.
\ce{XeF4} ($D_{4h}$): no. \ce{CH2Cl2} ($C_{2v}$): yes. cis-1,2-dichloroethene
($C_{2v}$): yes; trans ($C_{2h}$): no.
\end{solution}

\begin{solution}{exo:b3:point-groups:4}
$C_2(z) = \mathrm{diag}(-1,-1,1)$, $\chi = -1$; $i = \mathrm{diag}(-1,-1,-1)$,
$\chi = -3$; $\sigma_h(xy) = \mathrm{diag}(1,1,-1)$, $\chi = 1$.
\end{solution}

\begin{solution}{exo:b3:point-groups:5}
With the diagonal matrices of the previous exercise: $C_2i = \sigma_h$, $C_2\sigma_h =
i$, $i\sigma_h = C_2$, each operation is its own inverse, and every product
commutes. The group is commutative, so $X^{-1}AX = A$ for all $X$: each
operation is a class alone (four classes, four \omterm{def:b3:point-groups:irrep}{irreducible representations}
of dimension 1).
\end{solution}

\begin{solution}{exo:b3:point-groups:6}
Take $\sigma_v(1)$ the $xz$ plane. The point $(1,0)$ goes by $C_3$ to $(-\frac12,
\frac{\sqrt3}2)$, then by $\sigma_v(1)$ to $(-\frac12,-\frac{\sqrt3}2)$; in the other
order it goes to $(1,0)$ then $(-\frac12,\frac{\sqrt3}2)$. The two products are
reflections in two different planes ($\sigma_v(3)$ and $\sigma_v(2)$): the group is
not commutative, which is why its reflections form a class of three.
\end{solution}

\begin{solution}{exo:b3:point-groups:7}
Twisted biphenyl keeps the $C_2$ along the inter-ring bond and two $C_2$
perpendicular to it, bisecting the angles between the ring planes; every
plane and the \omterm{def:b3:point-groups:inversion}{centre of inversion} of the planar or perpendicular forms are
lost. With three perpendicular $C_2$ and no improper element the group is
$D_2$: the molecule is chiral (its two twisted forms are enantiomers,
separable when bulky ortho substituents prevent rotation).
\end{solution}

\begin{solution}{exo:b3:point-groups:8}
$d_{z^2}$ and $d_{x^2-y^2}$: $a_1$ ($x^2$, $y^2$, $z^2$ are $A_1$); $d_{xy}$: $a_2$; $d_{xz}$:
$b_1$; $d_{yz}$: $b_2$.
\end{solution}

\begin{solution}{exo:b3:point-groups:9}
$1 + 1 + 4 + 9 + 9 = 24 = h$. $\sum g\,\chi_{A_1}\chi_{T_2} = 1\cdot3 + 8\cdot0 + 3(-1)
+ 6(-1) + 6(1) = 0$.
\end{solution}

\begin{solution}{exo:b3:point-groups:10}
In the plane perpendicular to the intersection line, a reflection in a line
at angle $\alpha$ maps the polar angle $\phi$ to $2\alpha - \phi$. Two reflections, at
$\alpha$ then $\alpha + \theta$: $\phi \to 2\alpha - \phi \to 2(\alpha + \theta) - (2\alpha - \phi) =
\phi + 2\theta$, a rotation by $2\theta$ about the line. Two vertical planes at
$60^\circ$ thus generate a rotation by $120^\circ$: a $C_3$.
\end{solution}

\begin{solution}{exo:b3:point-groups:11}
The C=C=C axis is a $C_2$ and an $S_4$ (turn by $90^\circ$, which exchanges the
two \ce{CH2} planes, then reflect through the central carbon's plane); the two
\ce{CH2} planes are $\sigma_d$; two $C_2$ perpendicular to the axis bisect them.
Order 8: $D_{2d}$, achiral. In \ce{ClHC=C=CHCl} the planes, the $S_4$ and the axial
$C_2$ are destroyed by the substitution; only one $C_2$ perpendicular to the
axis survives: $C_2$, chiral --- an axially chiral allene.
\end{solution}

\begin{solution}{exo:b3:point-groups:12}
\ce{SF6}: $O_h$ ($h = 48$). \ce{SF5Cl}: the $C_4$ through Cl and four $\sigma_v$:
$C_{4v}$ ($h = 8$). trans-\ce{SF4Cl2}: $D_{4h}$ ($h = 16$). cis-\ce{SF4Cl2}: $C_{2v}$ ($h
= 4$). Each keeps a subset of the operations of $O_h$, closed under products,
and 8, 16, 4 divide 48. Polar: \ce{SF5Cl} and cis-\ce{SF4Cl2}.
\end{solution}

\begin{solution}{pb:b3:point-groups:1}
\textbf{1.} Along the four body diagonals through C and each H; each gives
$C_3$ and $C_3^2$: 8 rotations.
\textbf{2.} Along $x$, $y$, $z$ through the face centres: $C_2(z)$ maps $(1,1,1)
\to (-1,-1,1)$, a hydrogen. A rotation by $90^\circ$ about $z$ followed by
reflection in $xy$ maps $(1,1,1) \to (-1,1,1) \to (-1,1,-1)$, a hydrogen: an
$S_4$; with $S_4^3$, 6 operations.
\textbf{3.} The planes $x = \pm y$, $y = \pm z$, $x = \pm z$; $x = y$ contains
$(1,1,1)$ and $(-1,-1,1)$: each plane contains two C--H bonds.
\textbf{4.} $1 + 8 + 3 + 6 + 6 = 24$.
\textbf{5.} Inversion maps $(1,1,1)$ to $(-1,-1,-1)$, which is not a hydrogen:
no $i$.
\textbf{6.} $E$; $8C_3$; $3C_2$; $6S_4$; $6\sigma_d$.
\textbf{7.} \ce{CH3Cl}: $C_{3v}$, $h = 6$.
\textbf{8.} \ce{CH2Cl2}: $C_{2v}$, $h = 4$.
\textbf{9.} \ce{CHCl3}: $C_{3v}$; \ce{CCl4}: $T_d$, $h = 24$.
\textbf{10.} \ce{CH2FCl}: $C_s$ ($h = 2$); \ce{CHFClBr}: $C_1$ ($h = 1$).
\textbf{11.} Each keeps those operations of $T_d$ that respect the
substitution; 6, 4, 6, 24, 2, 1 all divide 24.
\textbf{12.} The three hydrogens, permuted by $C_3$, $C_3^2$ and the three
$\sigma_v$.
\textbf{13.} All but \ce{CCl4}: along the $C_3$ axis for \ce{CH3Cl} and \ce{CHCl3},
the $C_2$ axis for \ce{CH2Cl2}, in the \omterm{def:b3:point-groups:mirror}{mirror plane} for \ce{CH2FCl}, in no
imposed direction for \ce{CHFClBr}.
\textbf{14.} \ce{CHFClBr} only ($C_1$, no $S_n$).
\textbf{15.} Two enantiomers.
\textbf{16.} Its two hydrogens are equivalent: the plane through C, F and Cl
bisecting H--C--H is a \omterm{def:b3:point-groups:mirror}{mirror plane}.
\textbf{17.} No: four different substituents leave only $E$; with no $S_n$ the
molecule is chiral (and may be polar).
\textbf{18.} $(x, y, z)$: $T_2$.
\textbf{19.} Two hydrogens are the ends of an edge of the tetrahedron; the 24
operations carry any edge onto any other (the 6 edges form one set): one
\ce{CH2Cl2}.
\textbf{20.} \ce{CH2FCl}: one. \ce{CHFClBr}: two, a pair of enantiomers.
\textbf{21.} Three: ortho ($C_{2v}$), meta ($C_{2v}$), para ($D_{2h}$).
\textbf{22.} Three: 1,2,3- ($C_{2v}$), 1,2,4- ($C_s$), 1,3,5- ($D_{3h}$).
\textbf{23.} Two (Cl atoms adjacent or opposite). Only one \ce{CH2Cl2} is
known, which ruled out a planar carbon and supported the tetrahedral carbon
proposed in 1874.
\textbf{24.} $1^2 + 1^2 + 2^2 + 3^2 + 3^2 = 24$: \textbf{the order of $T_d$, the
\omterm{def:b3:point-groups:point-group}{point group} of methane, is $h = 24$}.
\end{solution}
